\documentclass[11pt,letterpaper]{article}
\usepackage[utf8]{inputenc}

%----- Configuración del estilo del documento------%
\usepackage{epsfig,graphicx}
\usepackage[left=2cm,right=2cm,top=1.8cm,bottom=2.3cm]{geometry}
\usepackage{fancyhdr}
\usepackage{lastpage}
\usepackage{url}
\pagestyle{fancy}
\fancyhf{}
\rfoot{\textit{Página \thepage \hspace{1pt} de \pageref{LastPage}}}


%------ Paquetes matemáticos básicos --------%
\usepackage{amsmath}
\usepackage{amssymb}
\usepackage{amsthm}

\usepackage[spanish]{babel}
\usepackage{graphicx}
\usepackage{hyperref}

\usepackage{tabularx}
\usepackage{xcolor}
\usepackage[table]{xcolor}
\usepackage{colortbl}
\usepackage{array, multirow, multicol, tabularx}
\usepackage{tcolorbox}
\newtheorem{theorem}{Theorem}[section]
\newtheorem{corollary}{Corollary}[theorem]
\newtheorem{lemma}[theorem]{Lemma}

%------si-------%
\definecolor{B}{HTML}{FFFFFF}
\definecolor{G}{HTML}{5e5e5e}
\definecolor{R2}{HTML}{d53d40}
\definecolor{A2}{HTML}{034190}
\definecolor{V2}{HTML}{7faa50}
\newcommand{\R}{\mathbb{R}}
\newcommand{\C}{\mathcal{C}}
\newcommand{\N}{\mathbb{N}}
\newcommand{\Z}{\mathbb{Z}}
\newcommand{\Q}{\mathbb{Q}}
\renewcommand{\theenumi}{\Roman{enumi}}
\renewcommand{\labelenumi}{{\theenumi}.}

\begin{document}

%------ Encabezado -------- %

\begin{center}
    \begin{minipage}{3cm}
    	\begin{center}
    		\includegraphics[height=3.4cm]{logo_unam.png}
    	\end{center}
    \end{minipage}\hfill
    \begin{minipage}{10cm}
    	\begin{center}
    	\textbf{\large Universidad Nacional Autónoma de México}\\[0.1cm]
        \textbf{Facultad de Ciencias}\\[0.1cm]
        \textbf{\'Algebra superior 2}\\[0.1cm]
        Tarea examen 1 \\[0.1cm]
         El\'ias L\'opez Rivera\\[0.1cm]
        \texttt{ elias.lopezr\,@ciencias.unam.mx }\\[0.1cm]
        Fecha:\,\,27/10/2024
    	\end{center}
    \end{minipage}\hfill
    \begin{minipage}{3cm}
    	\begin{center}
    		\includegraphics[height=3.4cm]{Logo_FC.png}
    	\end{center}
    \end{minipage}
\end{center}

\rule{17cm}{0.1mm}

%------ Fin de encabezado -------- %
\begin{tcolorbox}[
	title = \textcolor{black}{\textcolor{white}{Problema 1}},]
\textit{A partir de los naturales $\N$ define $\Z=(\N\times \N)/$ donde $(n,m) (k,m)$ si y solo si $n+l=m+k$
.Induce las operaciones de suma y producto en $\Z$ utilizando aquellas de $\N$, demuestra que est\'an bien definidas 
(en el sentido de que no dependen de los representantes) y demuestra que el resultado es un anillo conmutativo con $1$ que 
"contiene" a $\N$ 
}
\end{tcolorbox}\,\\
\begin{proof}
    
\end{proof}\,\\
\begin{tcolorbox}[
	title = \textcolor{black}{\textcolor{white}{Problema 2}},]
\textit{Demuestra que para todo elemento $[(n,m)]\in \Z$ se satisface exactamente una de las siguientes afirmaciones:
\begin{itemize}
    \item Existe $c\in \N$ tal que $[(n,m)]=[(c,0)]$
    \item $[(n,m)]=[(0,0)]$
    \item Existe $c\in \N$ tal que $[(n,m)]=[(0,c)]$
\end{itemize}
}
\end{tcolorbox}\,\\
\begin{proof}
    
\end{proof}\,\\
\begin{tcolorbox}[
	title = \textcolor{black}{\textcolor{white}{Problema 3}},]
\textit{Demuestra que un elemento $u$ de un dominio entero tiene a lo m\'as un inverso multiplicativo
}
\end{tcolorbox}\,\\
\begin{proof}
    
\end{proof}\,\\
\begin{tcolorbox}[
	title = \textcolor{black}{\textcolor{white}{Problema 4}},]
\textit{Demuestra:\,$\forall\,\,a,b\in \Z\,(a>0\,\,b>0\,\,a^2>b^2 \implies a>b)$
}
\end{tcolorbox}\,\\
\begin{proof}\,\\

    
\end{proof}\,\\
\begin{tcolorbox}[
	title = \textcolor{black}{\textcolor{white}{Problema 5}},]
\textit{Demuestra que un dominio entero $R$ se vale la ley de cancelaci\'on:\,\,$\forall\,\,a,b,c\in R\,(a\neq 0\,\,ab=ac \implies b=c)$
}
\end{tcolorbox}\,\\
\begin{proof}\,\\
    \,\\
    Tomamos  $a,b,c\in \Z$ tal que $a\neq 0$, tenemos que:\,\\
    \,\\
    \begin{equation*}
        ab=ac
    \end{equation*}\,\\
    Como $ac$ es un elemento del dominio entero $R$, entonces este tiene un \'unico inverso $-(ac)\in R$, sumando por ambos
    lados tenemos que:\,\\
    \begin{equation*}
        ab+(-ac)=ac+(-ac)
    \end{equation*}\,\\
    Como $-(ac)$ se definio como el inverso bajo $+$ de $ac$, se sigue que:\,\\
    \begin{equation*}
        ab+(-ac)=0
    \end{equation*}\,\\
    Como $ab,-(ac)\in R$, podemos aplicar la distribuci\'on del producto sobre la suma:\,\\
    \begin{equation*}
        a(b+(-c))=0
    \end{equation*}\,\\
    Como $R$ es dominio entero necesariamente $a=0$ o $(b+(-c))=0$, como por hip\'otesis $a\neq 0$, entonces
    $(b+(-c))=0$, por tanto sumando $c$ a ambos lados de la igualdad:\,\\
    \begin{equation*}
        b+(-c)+c=0+c\implies b+0=0+c\implies b=c
    \end{equation*}
\end{proof}\,\\
\begin{tcolorbox}[
	title = \textcolor{black}{\textcolor{white}{Problema 1}},]
\textit{Demuestra que un elemento $\displaystyle \sum_{n=0}^{\infty}\,a_n\,x^n$ del anillo de series formales $\Q\,[[x]]$
con coeficientes en $\Q$ es una unidad en $\Q[[x]]$ si y solo si $a_0\neq 0$
}
\end{tcolorbox}\,\\
\begin{proof}
    
\end{proof}\,\\
\end{document}